% from ega0-1-fr.tex
\oldpage[0\textsubscript{I}]{52}

la bijection ainsi définie est un homomorphisme de groupe. Or, étant donnés
deux $\mathcal{O}_X$-Modules inversibles $\mathcal{L}$, $\mathcal{L}'$, il y a
un recouvrement ouvert $(U_\lambda)$ tel que $\mathcal{L}|U_\lambda$ et
$\mathcal{L}'|U_\lambda$ soient isomorphes à $\mathcal{O}_X|U_\lambda$ pour
tout $\lambda$~; il y a donc pour chaque indice $\lambda$ un élément
$a_\lambda$ (resp. $a'_\lambda$) de $\Gamma(U_\lambda,\mathcal{L})$ (resp.
$\Gamma(U_\lambda,\mathcal{L}'))$ tel que les éléments de
$\Gamma(U_\lambda,\mathcal{L})$ (resp. $\Gamma(U_\lambda,\mathcal{L}'))$
soient les $s_\lambda\cdot a_\lambda$ (resp.
$s_\lambda\cdot a'_\lambda$), où $s_\lambda$ parcourt
$\Gamma(U_\lambda,\mathcal{O}_X)$. Les cocycles correspondants
$(\varepsilon_{\lambda\mu})$, $(\varepsilon'_{\lambda\mu})$ sont tels que
$s_\lambda\cdot a_\lambda=s_\mu\cdot a_\mu$ (resp.
$s_\lambda\cdot a'_\lambda=s_\mu\cdot a'_\mu$) au-dessus de
$U_\lambda\cap U_\mu$ soit équivalent à
$s_\lambda=\varepsilon_{\lambda\mu}s_\mu$ (resp.
$s_\lambda=\varepsilon'_{\lambda\mu}s_\mu$) au-dessus de
$U_\lambda\cap U_\mu$. Comme les sections de
$\mathcal{L}\otimes_{\mathcal{O}_X}\mathcal{L}'$ au-dessus de $U_\lambda$
sont les sommes finies des
$s_\lambda s'_\lambda\cdot(a_\lambda\otimes a'_\lambda)$ où $s_\lambda$ et
$s'_\lambda$ parcourent $\Gamma(U_\lambda,\mathcal{O}_X)$, il est clair que
le cocycle $(\varepsilon_{\lambda\mu}\varepsilon'_{\lambda\mu})$ correspond
à $\mathcal{L}\otimes_{\mathcal{O}_X}\mathcal{L}'$, ce qui achève la
démonstration\footnote{Pour une forme générale de ce résultat, voir le livre
cité dans la note de la p.~51.}.
\end{env}

\begin{env}[5.4.8]
\label{0.5.4.8-fr}
Soit $f=(\psi,\omega)$ un morphisme $Y\to X$ d'espaces annelés. Le foncteur
$f^*(\mathcal{L})$ dans la catégorie des $\mathcal{O}_X$-Modules libres
définit une application (encore notée $f^*$ par abus de langage) de l'ensemble
$\mathfrak{M}(X)$ dans l'ensemble $\mathfrak{M}(Y)$. D'autre part, on a un
homomorphisme canonique (T, 3.2.2)
\[
  H^1(X,\mathcal{O}_X^*)\longrightarrow H^1(Y,\mathcal{O}_Y^*).
  \tag{5.4.8.1}\label{0.5.4.8.1-fr}
\]
Lorsqu'on identifie canoniquement (5.4.7) $\mathfrak{M}(X)$ et
$H^1(X,\mathcal{O}_X^*)$ (resp. $\mathfrak{M}(Y)$ et
$H^1(Y,\mathcal{O}_Y^*)$), l'homomorphisme (5.4.8.1) \emph{s'identifie à
l'application $f^*$}. En effet, si $\mathcal{L}$ provient d'un cocycle
$(\varepsilon_{\lambda\mu})$ correspondant à un recouvrement ouvert
$(U_\lambda)$ de $X$, il suffit de montrer que $f^*(\mathcal{L})$ provient
d'un cocycle dont la classe de cohomologie est image par (5.4.8.1) de celle de
$(\varepsilon_{\lambda\mu})$. Or, si $\theta_{\lambda\mu}$ est
l'automorphisme de $\mathcal{O}_X|(U_\lambda\cap U_\mu)$ qui correspond à
$\varepsilon_{\lambda\mu}$, il est clair que $f^*(\mathcal{L})$ s'obtient par
recollement des $\mathcal{O}_Y|\psi^{-1}(U_\lambda)$ au moyen des
automorphismes $f^*(\theta_{\lambda\mu})$, et il suffit donc de vérifier que
ces derniers correspondent au cocycle
$(\omega^\#(\varepsilon_{\lambda\mu}))$, ce qui résulte aussitôt des
définitions (on a identifié $\varepsilon_{\lambda\mu}$ à son image canonique
par $\rho$ (3.7.2), section de $\psi^*(\mathcal{O}_X^*)$ au-dessus de
$\psi^{-1}(U_\lambda\cap U_\mu)$).
\end{env}

\begin{env}[5.4.9]
\label{0.5.4.9-fr}
Soient $\mathcal{E}$, $\mathcal{F}$ deux $\mathcal{O}_X$-Modules,
$\mathcal{F}$ étant supposé \emph{localement libre}, et soit $\mathcal{G}$ un
$\mathcal{O}_X$-Module \emph{extension de $\mathcal{F}$ par $\mathcal{E}$},
autrement dit tel qu'il existe une suite exacte
\[
  0\longrightarrow\mathcal{E}\xrightarrow{i}\mathcal{G}
  \xrightarrow{p}\mathcal{F}\longrightarrow 0.
\]
Alors, pour tout $x\in X$, il existe un voisinage ouvert $U$ de $x$ tel que
$\mathcal{G}|U$ soit isomorphe à la \emph{somme directe}
$\mathcal{E}|U\oplus\mathcal{F}|U$. On peut en effet se borner au cas où
$\mathcal{F}=\mathcal{O}_X^n$~; soient $e_i$ ($1\leq i\leq n$) les sections
canoniques (5.5.5) de $\mathcal{O}_X^n$~; il existe alors un voisinage ouvert
$U$ de $x$ et $n$ sections $s_i$ de $\mathcal{G}$ au-dessus de $U$ telles que
$p(s_i|U)=e_i|U$ pour $1\leq i\leq n$. Cela étant, soit $f$
l'homomorphisme $\mathcal{F}|U\to\mathcal{G}|U$ défini par les sections
$s_i|U$ (5.1.1). Il est immédiat que pour tout ouvert $V\subset U$, et toute
section $s\in\Gamma(V,\mathcal{G})$ on a
$s-f(p(s))\in\Gamma(V,\mathcal{E})$, d'où notre assertion.
\end{env}

\begin{env}[5.4.10]
\label{0.5.4.10-fr}
Soient $f:X\to Y$ un morphisme d'espaces annelés, $\mathcal{F}$ un
$\mathcal{O}_X$-Module, $\mathcal{L}$ un $\mathcal{O}_Y$-Module localement
libre de rang fini. Il existe alors un isomorphisme canonique
\[
  f_*(\mathcal{F})\otimes_{\mathcal{O}_Y}\mathcal{L}
  \xrightarrow{\sim}
  f_*\bigl(\mathcal{F}\otimes_{\mathcal{O}_X}f^*(\mathcal{L})\bigr).
  \tag{5.4.10.1}\label{0.5.4.10.1-fr}
\]

